\chapter{Đào sâu: Trình duyệt gặp gateway --- nơi cất token, luồng refresh và CORS}
\label{ch:dd-frontend-auth}

Chương~26 đặt ứng dụng React lên cùng origin với API và ghi chú, gần
như thoáng qua, rằng CORS ``đơn giản là không bao giờ được hỏi tới''.
Chương~28 đẩy thiết kế token từ phía server. Chương đào sâu này lấy đầu
kia của đường dây: trình duyệt, runtime duy nhất trong hệ thống mà bạn
không kiểm soát, nơi chạy code do bất kỳ ai chèn được thẻ script vào
trang của bạn, và sẵn lòng gửi cookie tới form của kẻ tấn công. Mọi
quyết định về nơi một token sống là quyết định về mối đe dọa nào trong
số đó bạn chấp nhận.

Ba câu hỏi sắp xếp chương này:

\begin{enumerate}
  \item \textbf{Token sống ở đâu} trong trình duyệt, và mỗi chỗ ở
  phơi nó ra trước đòn tấn công nào?
  \item \textbf{Chuyện gì xảy ra khi hết hạn} --- luồng refresh chạy
  thế nào, và vì sao mười tab phải sinh ra đúng một lời gọi refresh?
  \item \textbf{CORS bảo vệ cái gì}, nó tốn gì mỗi request, và vì sao
  một origin làm nó biến mất?
\end{enumerate}

\section{Vì sao chọn nó, và nó giải quyết gì}

Server phát hai thông tin xác thực (Chương~5): một JWT access token 15
phút và một refresh token mờ 7 ngày. Trình duyệt phải giữ cả hai qua
các lần tải lại trang, đính cái thứ nhất vào mọi lời gọi API, và đổi
cái thứ hai lấy cái thứ nhất mới khi hết hạn --- mà người dùng không
nhận ra. Câu trả lời của dự án là một store Zustand với middleware
\code{persist} và một interceptor của axios. Đọc store trước:

\begin{lstlisting}
export const useAuthStore = create<AuthState>()(
  persist(                                            (*@\co{1}@*)
    (set) => ({
      accessToken: null,
      refreshToken: null,
      user: null,
      isAuthenticated: false,
      setTokens: (accessToken, refreshToken) =>
        set({ accessToken, refreshToken, isAuthenticated: true }),
      setUser: (user) => set({ user }),
      logout: () =>
        set({ accessToken: null, refreshToken: null,
              user: null, isAuthenticated: false }),
    }),
    { name: 'auth-storage' }                          (*@\co{2}@*)
  )
);
\end{lstlisting}

\begin{description}
  \item[\co{1}] \code{persist} ghi toàn bộ state vào storage ở mỗi
  thay đổi và nạp lại khi tải trang. Đó là thứ sống sót qua một lần
  tải lại; không có nó, mỗi lần làm mới trang là một lần đăng xuất.
  \item[\co{2}] Khóa lưu trữ. Backend mặc định là \code{localStorage},
  nên mở DevTools, Application, Local Storage, và chúng ở đó: cả hai
  token, dạng chữ thường, dưới \code{auth-storage}. Bất kỳ script nào
  chạy trên trang đọc được chúng bằng một dòng.
\end{description}

\begin{decision}[localStorage, cho phiên bản đầu]
Dự án cất cả hai token trong \code{localStorage} vì đó là thứ đơn giản
nhất sống sót qua tải lại trang, hoạt động giống hệt trên
\code{localhost:3000} và \code{ts.local}, không cần cấu hình cookie ở
server, và né hẳn CSRF --- một token do JavaScript tự tay đính vào
không bao giờ được trình duyệt tự gửi. Cái giá là token chỉ an toàn
đúng bằng mọi script trên trang. Điều đó chấp nhận được khi trang chỉ
tải bundle của chính nó. Nó hết chấp nhận được vào ngày một script bên
thứ ba (một thẻ analytics, một widget chat, một dependency npm bị
chiếm) chạy trên cùng origin: bất kỳ cái nào cũng đọc được
\code{auth-storage} và chuyển một refresh token 7 ngày đi nơi khác.
Nâng cấp ở Mục~\ref{sec:fe-enhance} chuyển refresh token vào cookie
\code{HttpOnly} và giữ access token trong bộ nhớ --- nhiều bộ phận
chuyển động hơn, nhưng script không còn trộm được thông tin xác thực
sống dài.
\end{decision}

\section{Luồng refresh, như đã viết}

\subsection{Interceptor}

Mọi lời gọi API đi qua một instance axios:

\begin{lstlisting}
const apiClient = axios.create({ baseURL: '/api' });   (*@\co{1}@*)

apiClient.interceptors.request.use((config) => {
  const token = useAuthStore.getState().accessToken;
  if (token) config.headers.Authorization = `Bearer ${token}`;
  return config;                                        (*@\co{2}@*)
});

apiClient.interceptors.response.use(
  (response) => response,
  async (error) => {
    const originalRequest = error.config;
    if (error.response?.status === 401
        && !originalRequest._retry) {                   (*@\co{3}@*)
      originalRequest._retry = true;
      const refreshToken = useAuthStore.getState().refreshToken;
      if (refreshToken) {
        try {
          const { data } = await axios.post(            (*@\co{4}@*)
              '/api/auth/refresh', { refreshToken });
          useAuthStore.getState()
              .setTokens(data.accessToken, data.refreshToken);
          originalRequest.headers.Authorization =
              `Bearer ${data.accessToken}`;
          return apiClient(originalRequest);            (*@\co{5}@*)
        } catch {
          useAuthStore.getState().logout();
          window.location.href = '/login';
        }
      }
    }
    return Promise.reject(error);
  }
);
\end{lstlisting}

\begin{description}
  \item[\co{1}] Base URL tương đối. Khi phát triển, proxy của Vite
  chuyển \code{/api} tới \code{localhost:8080}; trong cụm Traefik làm
  việc đó (Chương~26). Trình duyệt không bao giờ thấy origin thứ hai.
  \item[\co{2}] Header bearer do code đính vào mọi request. Đây là sự
  miễn nhiễm CSRF nhắc ở trên: một form trên \code{evil.example} không
  thể khiến trình duyệt thêm header này.
  \item[\co{3}] Một 401 từ gateway kích hoạt một lần thử lại cho mỗi
  request; body JSON ``Invalid or expired token'' đến từ filter
  \code{JwtAuthenticationFilter} của gateway. Cờ được đánh dấu trên
  object config để 401 thứ hai rơi xuống bên gọi.
  \item[\co{4}] Lời gọi refresh cố ý dùng \code{axios} \emph{trần},
  không phải \code{apiClient}: nếu không, một 401 từ chính
  \code{/refresh} sẽ vào lại interceptor này và lặp.
  \item[\co{5}] Request gốc được phát lại với token mới. \code{await}
  của bên gọi hoàn tất như thể không có gì xảy ra.
\end{description}

Điều này đúng cho một request. Nó sai cho một trang.

\subsection{Bài toán single-flight}

\code{DashboardPage} mount và TanStack Query bắn các truy vấn song
song: khóa học, bài tập, hồ sơ người dùng, số đếm từ điển. Nếu access
token hết hạn khi laptop đóng nắp, tất cả trả 401 trong cùng vài mili
giây, và mỗi cái chạy interceptor độc lập. Bốn request, bốn lời gọi tới
\code{/api/auth/refresh}, bốn access token mới, và --- vì
\code{setTokens} chạy bốn lần --- phản hồi nào tới sau cùng thì thắng.
Hôm nay điều đó vô hại vì auth-service trả cùng một refresh token mỗi
lần (Chương~28 đã chỉ ra đây là việc xoay vòng còn thiếu). Thêm xoay
vòng, điều bạn nên làm, và cuộc đua thành chí mạng: lần refresh đầu vô
hiệu token mà ba cái kia sắp đưa ra, chúng thất bại, và người dùng bị
đăng xuất bởi chính cơ chế lẽ ra giữ họ lại.

Cách sửa là refresh \emph{single-flight}: 401 đầu tiên khởi động một
lần refresh và lưu promise của nó; mọi 401 đồng thời await cùng
promise đó thay vì tự khởi động; khi nó hoàn tất, tất cả phát lại với
cùng token mới. Mục~\ref{sec:fe-enhance} có code đầy đủ. Người phỏng
vấn hỏi câu này dưới dạng ``mười tab gặp 401 cùng lúc --- bao nhiêu lời
gọi refresh?'' và câu trả lời senior là ``mỗi tab một nếu không phối
hợp, tổng cộng một nếu có khóa xuyên tab'', theo sau là gọi tên công cụ
cho trường hợp xuyên tab: \code{Web Locks} API hoặc một
\code{BroadcastChannel}, vì mỗi tab có heap JavaScript riêng và không
thấy promise đang chờ của tab khác.

\subsection{Đăng xuất và trang callback OAuth}

\code{Header.handleLogout} POST refresh token tới
\code{/api/auth/logout}, nơi xóa dòng của nó, rồi dọn store và điều
hướng tới \code{/login}. Việc đó thu hồi thông tin xác thực 7 ngày ngay
lập tức. Access token 15 phút không bị thu hồi --- không thể, theo
Chương~28 --- nên một bản sao lấy trước khi đăng xuất vẫn hoạt động cho
tới khi hết hạn. ``Làm sao đăng xuất một người dùng ở mọi nơi?'' có
cùng câu trả lời với thêm một bước: xóa \emph{tất cả} dòng refresh
token của người dùng, không chỉ dòng trình duyệt này giữ, và chấp nhận
cửa sổ 15 phút hoặc thêm deny-list theo \code{jti}.

Đăng nhập Google hạ cánh ở \code{OAuthCallbackPage} với cả hai token
là tham số truy vấn --- \code{OAuth2LoginSuccessHandler} dựng redirect
theo cách đó. Trang đọc chúng, cất chúng, và điều hướng với
\code{replace: true} để URL chứa token không nằm lại trong lịch sử.
Chi tiết cuối quan trọng hơn vẻ ngoài: một URL được chép vào lịch sử
trình duyệt, vào header \code{Referer} của bất kỳ tài nguyên nào trang
tải tiếp theo, và vào access log của server. \code{replace} rút ngắn
thời gian phơi nhiễm; nó không xóa được khỏi access log của gateway
và auth-service, nơi giờ chứa một refresh token hợp lệ dạng chữ
thường.

\section{CORS: chính sách bạn không đang dùng}

Chính sách same-origin là quy tắc của trình duyệt rằng một script tải
từ origin A không được \emph{đọc} phản hồi từ origin B. CORS là giao
thức để B nới lỏng quy tắc đó theo từng origin. Gateway của dự án có
một \code{CorsConfig} liệt kê \code{localhost:3000} và
\code{localhost:5173}, cho phép mọi method và header, và đặt
\code{allowCredentials(true)}. Đọc nó đối chiếu với hai cách ứng dụng
thực sự được phục vụ:

\begin{itemize}
  \item \textbf{Phát triển.} Vite phục vụ trang ở \code{:3000} và
  \emph{proxy} \code{/api} tới \code{:8080}. Trình duyệt gửi request
  tới \code{:3000}; Vite chuyển tiếp phía server. Không tồn tại request
  cross-origin, nên không preflight, nên \code{CorsConfig} không bao
  giờ được dùng.
  \item \textbf{Cụm.} Traefik định tuyến \code{/api} tới gateway và
  mọi thứ khác tới nginx trên cùng host (Chương~26). Lại cùng origin.
\end{itemize}

\code{CorsConfig} vì thế phục vụ đúng một kịch bản: ai đó chạy bundle
đã build từ một origin khác --- chẳng hạn mở \code{dist/index.html}
qua một static server ở \code{:5173} không có proxy. Nó là lưới an
toàn, và là lưới rộng rãi: \code{addAllowedHeader("*")} cộng
credentials là chính sách rộng nhất Spring chấp nhận mà không ném lỗi.

CORS sẽ tốn gì nếu nó hoạt động là điều đáng biết cho phỏng vấn. Một
request với body JSON và header \code{Authorization} không phải
``request đơn giản'', nên trình duyệt gửi trước một preflight
\code{OPTIONS} hỏi method và header có được phép không. Đó là thêm một
vòng đi-về mỗi request khác nhau, có thể cache bằng
\code{Access-Control-Max-Age}, thứ cấu hình của dự án không đặt --- nên
mọi lời gọi sẽ preflight. Trên đường truyền 40\,ms, mọi lời gọi API
thành 80\,ms trước khi có việc gì được làm. Một origin loại bỏ hoàn
toàn chi phí đó, và đó là lý lẽ hiệu năng cho thiết kế của Chương~26
bên cạnh lý lẽ bảo mật.

\section{Ưu và nhược}

\begin{center}
\begin{tabular}{@{}p{0.30\linewidth}p{0.32\linewidth}p{0.32\linewidth}@{}}
\toprule
 & \textbf{localStorage + bearer (hôm nay)} & \textbf{Cookie HttpOnly} \\
\midrule
XSS trộm token & có, cả hai token, một dòng JS & không --- script không
đọc được cookie \\
CSRF & miễn nhiễm: header do code đính & phơi nhiễm trừ khi có
\code{SameSite} + token CSRF hoặc kiểm tra header tùy biến \\
Sống qua tải lại & có, qua \code{persist} & có, trình duyệt giữ nó \\
Nhiều tab & mỗi tab đọc storage; refresh đua nhau & cùng cookie ở mọi
tab; refresh vẫn đua \\
Chạy cross-origin & có & cần \code{Allow-Credentials} và danh sách
origin chính xác \\
Thay đổi server & không & set-cookie khi đăng nhập, đọc cookie khi
refresh \\
\bottomrule
\end{tabular}
\end{center}

Tóm tắt thành thật mà một kỹ sư senior đưa ra: bearer-trong-storage
đổi một bài toán CSRF bạn giải được lấy một bài toán XSS bạn không giải
được. Bất kỳ script bị chèn nào đánh bại nó hoàn toàn. Cookie đổi
theo chiều ngược lại, và bài toán của nó có cách sửa cơ học đã biết.

\section{Cái gì gãy}

\begin{pitfall}[bốn lần refresh, một người thắng]
Triệu chứng: sau khi để ứng dụng mở qua đêm, cú click đầu tiên chạy,
rồi người dùng bị đưa tới \code{/login} dù tài khoản không có gì sai.
Nguyên nhân: các 401 song song mỗi cái chạy interceptor; với xoay vòng
refresh token ở server, lần refresh thứ hai đưa ra token mà lần đầu
vừa vô hiệu, và nhánh \code{catch} đăng xuất người dùng. Cách sửa:
interceptor single-flight ở Mục~\ref{sec:fe-enhance}. Không có xoay
vòng, triệu chứng chỉ là lời gọi lãng phí --- đó là lý do cuộc đua vô
hình hôm nay.
\end{pitfall}

\begin{pitfall}[token trong URL]
Triệu chứng: một refresh token xuất hiện trong access log của gateway,
trong log của Traefik, và trong lịch sử trình duyệt của một máy dùng
chung. Nguyên nhân: \code{OAuth2LoginSuccessHandler} truyền cả hai
token làm tham số truy vấn của redirect tới \code{/oauth/callback}.
Cách sửa: đặt token vào \emph{fragment} của URL
(\code{\#accessToken=...}) --- fragment không bao giờ được gửi tới
server nào --- hoặc tốt hơn, chỉ trả về một mã dùng một lần mà trang
đổi lấy token bằng POST. Thiết kế cookie bên dưới loại refresh token
khỏi URL hoàn toàn.
\end{pitfall}

\begin{pitfall}[một 401 không liên quan tới token]
Triệu chứng: mật khẩu sai ở trang đăng nhập làm khung ứng dụng lóe lên
chốc lát, rồi quay về \code{/login}. Nguyên nhân: auth-service trả lời
đăng nhập sai bằng 401; interceptor coi \emph{mọi} 401 là ``token hết
hạn'' và, nếu một refresh token cũ còn trong storage từ phiên trước,
thử refresh, thất bại, rồi chạy \code{logout()} cộng chuyển hướng
cứng. Cách sửa: miễn các endpoint auth khỏi đường refresh
(\code{originalRequest.url} bắt đầu bằng \code{/auth/}), và dọn store
khi trang đăng nhập mount.
\end{pitfall}

\section{Chuyện gì gãy lúc 3 giờ sáng}

Ba sự cố chỉ client thấy.

\begin{itemize}
  \item \textbf{Laptop đóng nắp tám ngày.} Khi quay lại, access token
  đã chết từ lâu và refresh token hết hạn hôm qua. Request đầu nhận
  401, refresh nhận 401 (auth-service xóa dòng hết hạn và ném lỗi),
  interceptor đăng xuất người dùng. Hành vi đúng --- nhưng mọi truy
  vấn đang bay trên trang thất bại trước, và TanStack Query có thể thử
  lại mỗi cái ba lần trước khi chuyển hướng tới nơi. Đặt
  \code{retry: false} cho 401 trong tùy chọn mặc định của query client,
  nếu không người dùng nhìn vòng xoay mười giây trước khi được đưa tới
  đăng nhập.
  \item \textbf{Đồng hồ client lệch.} Dự án không bao giờ kiểm tra
  \code{exp} ở client, và thế là đúng: đồng hồ client không đáng tin
  và gateway quyết định. Một client ``thông minh'' chủ động refresh khi
  \code{exp} sắp tới sẽ, trên máy có đồng hồ nhanh 20 phút, refresh ở
  mọi request --- hoặc, chậm 20 phút, không bao giờ refresh và vẫn
  nhận 401. Hãy để 401 của server là tín hiệu.
  \item \textbf{CORS cấu hình sai chỉ ở production.} Triệu chứng là lỗi
  console ``No Access-Control-Allow-Origin header'' trên một request
  trả 200 trong tab network. Trình duyệt nhận được phản hồi và từ chối
  trao nó cho script. Nó xảy ra khi trang và API ở hai origin khác nhau
  trong production (CDN cho bundle, gateway ở host khác) và danh sách
  cho phép của gateway vẫn ghi \code{localhost}. Cách kiểm tra luôn
  giống nhau: so header request \code{Origin} với danh sách
  \code{allowedOrigins}, từng ký tự, cả scheme và cổng. Và để ý hình
  dạng của lỗi: nó không bao giờ hiện với server, nơi log ghi một
  request thành công.
\end{itemize}

\section{Nâng cấp giải pháp}
\label{sec:fe-enhance}

Thiết kế đích: access token chỉ sống trong bộ nhớ JavaScript (mất khi
tải lại, khôi phục bằng một lần refresh âm thầm); refresh token sống
trong cookie \code{HttpOnly} mà JavaScript không đọc được và trình
duyệt chỉ gửi tới endpoint refresh.

\textbf{1. Auth-service đặt cookie.} Khi đăng nhập, xác minh 2FA, và
refresh, controller trả access token trong body và refresh token trong
cookie:

\begin{lstlisting}[language=Java]
private ResponseCookie refreshCookie(String value, long maxAgeSec) {
    return ResponseCookie.from("refresh_token", value)
            .httpOnly(true)       // no document.cookie access
            .secure(true)         // HTTPS only; false on plain localhost
            .sameSite("Strict")   // never sent from another site
            .path("/api/auth/refresh")   // only this endpoint sees it
            .maxAge(maxAgeSec)
            .build();
}

@PostMapping("/login")
public ResponseEntity<LoginResponse> login(
        @Valid @RequestBody LoginRequest request) {
    LoginResponse r = authService.login(request);
    if (r.refreshToken() == null) {          // 2FA or password change
        return ResponseEntity.ok(r);
    }
    return ResponseEntity.ok()
            .header(HttpHeaders.SET_COOKIE, refreshCookie(
                    r.refreshToken(),
                    jwtService.getRefreshTokenExpiryMs() / 1000).toString())
            .body(new LoginResponse(r.accessToken(), null,
                    false, false, null));    // token not in the body
}

@PostMapping("/refresh")
public ResponseEntity<LoginResponse> refresh(
        @CookieValue(name = "refresh_token", required = false)
        String cookie) {
    if (cookie == null) {
        return ResponseEntity.status(HttpStatus.UNAUTHORIZED).build();
    }
    LoginResponse r = authService.refreshToken(
            new TokenRefreshRequest(cookie));
    return ResponseEntity.ok()
            .header(HttpHeaders.SET_COOKIE, refreshCookie(
                    r.refreshToken(),
                    jwtService.getRefreshTokenExpiryMs() / 1000).toString())
            .body(new LoginResponse(r.accessToken(), null,
                    false, false, null));
}

@PostMapping("/logout")
public ResponseEntity<Void> logout(
        @CookieValue(name = "refresh_token", required = false)
        String cookie) {
    if (cookie != null) authService.logout(cookie);
    return ResponseEntity.noContent()
            .header(HttpHeaders.SET_COOKIE,
                    refreshCookie("", 0).toString())   // delete it
            .build();
}
\end{lstlisting}

\code{path} của cookie là thuộc tính mạnh nhất trong bốn: ngay cả khi
một endpoint nào khác có lỗ hổng, trình duyệt không bao giờ đính cookie
này vào nó. Để ý path là góc nhìn của \emph{gateway}
(\code{/api/auth/refresh}), vì trình duyệt nói chuyện với gateway;
gateway bỏ tiền tố sau đó. \code{SameSite=Strict} là phòng thủ CSRF:
một form POST xuyên site không mang cookie chút nào. Yêu cầu body JSON
trên \code{/refresh} là lớp thứ hai --- một form HTML thuần không gửi
được \code{application/json} mà không có preflight, thứ thiết kế
same-origin sẽ không cấp cho site khác.

\textbf{2. Client giữ access token trong bộ nhớ.} Store bỏ
\code{persist} và refresh token hoàn toàn:

\begin{lstlisting}
export const useAuthStore = create<AuthState>()((set) => ({
  accessToken: null,          // memory only: gone on reload, by design
  user: null,
  isAuthenticated: false,
  setAccessToken: (accessToken) =>
    set({ accessToken, isAuthenticated: true }),
  setUser: (user) => set({ user }),
  logout: () =>
    set({ accessToken: null, user: null, isAuthenticated: false }),
}));
\end{lstlisting}

Khi tải trang, bộ nhớ trống; ứng dụng thực hiện một lần refresh âm
thầm trước khi render các route được bảo vệ. Một lần tải lại vì thế tốn
một vòng đi-về tới \code{/api/auth/refresh}, và nếu cookie không còn,
người dùng ở \code{/login} --- cùng hành vi như phiên hết hạn.

\textbf{3. Refresh single-flight với phát lại.} Interceptor giữ một
promise đang chờ và một hàng đợi request chờ nó:

\begin{lstlisting}
let refreshing: Promise<string> | null = null;   // module-level lock

async function refreshAccessToken(): Promise<string> {
  if (!refreshing) {                              // first 401 starts it
    refreshing = axios
      .post('/api/auth/refresh', null, { withCredentials: true })
      .then(({ data }) => {
        useAuthStore.getState().setAccessToken(data.accessToken);
        return data.accessToken as string;
      })
      .finally(() => { refreshing = null; });     // release the lock
  }
  return refreshing;                              // others await it
}

apiClient.interceptors.response.use(
  (response) => response,
  async (error) => {
    const original = error.config;
    const isAuthCall = original.url?.startsWith('/auth/');
    if (error.response?.status !== 401 || original._retry
        || isAuthCall) {
      return Promise.reject(error);
    }
    original._retry = true;
    try {
      const token = await refreshAccessToken();    // shared promise
      original.headers.Authorization = `Bearer ${token}`;
      return apiClient(original);                  // replay
    } catch {
      useAuthStore.getState().logout();
      window.location.href = '/login';
      return Promise.reject(error);
    }
  }
);
\end{lstlisting}

Mười 401 đồng thời giờ sinh ra một POST; chín cái kia \code{await}
cùng promise và phát lại với token nó trả về. \code{withCredentials:
true} là thứ khiến axios gửi cookie --- và là thứ khiến request thành
\emph{credentialed} với CORS. Để phối hợp xuyên tab, bọc thân của
\code{refreshAccessToken} trong \code{navigator.locks.request('refresh',
...)}: Web Locks API là toàn tiến trình, nên các tab thay phiên, và lần
refresh của tab thứ hai thấy cookie đã được xoay bởi tab đầu và đơn
giản là thành công.

\textbf{4. Gateway, nếu có ngày xuất hiện origin thứ hai.} Request
credentialed cấm wildcard; origin phải được liệt kê chính xác và
preflight được cache:

\begin{lstlisting}[language=yaml]
spring:
  cloud:
    gateway:
      globalcors:
        cors-configurations:
          '[/**]':
            allowed-origins:
              - https://app.teachersupporter.example   # exact, with scheme
            allowed-methods: [GET, POST, PUT, PATCH, DELETE, OPTIONS]
            allowed-headers: [Authorization, Content-Type]
            allow-credentials: true                   # cookies cross-origin
            max-age: 3600                             # cache the preflight
\end{lstlisting}

Với \code{SameSite=Strict} cookie vẫn không đi tới một site khác, nên
một triển khai thực sự xuyên site cần \code{SameSite=None; Secure} và
chấp nhận phơi nhiễm CSRF đi kèm --- lý do thiết kế same-origin của
Chương~26 là thứ nên giữ.

\begin{handson}[xem cookie làm việc của nó]
Sau khi triển khai phiên bản cookie, đăng nhập qua UI, rồi kiểm tra
trình duyệt giữ gì và gửi gì:
\begin{lstlisting}[language=bash]
$ curl -si -c jar.txt -X POST http://localhost:8080/api/auth/login \
    -H 'Content-Type: application/json' \
    -d '{"email":"teacher@ts.local","password":"secret"}' | grep -i set-cookie
Set-Cookie: refresh_token=3f9c...; Path=/api/auth/refresh; Max-Age=604800;
  HttpOnly; SameSite=Strict
$ curl -s -b jar.txt -X POST http://localhost:8080/api/auth/refresh \
    | jq .accessToken
"eyJhbGciOiJIUzI1NiJ9..."
$ curl -s -b jar.txt -o /dev/null -w '%{http_code}\n' \
    http://localhost:8080/api/courses
401
\end{lstlisting}
Dòng cuối là thuộc tính \code{path} đang làm việc: cookie tồn tại
trong jar nhưng curl, giống trình duyệt, không gửi nó tới
\code{/api/courses}. Trong DevTools cùng cookie đó hiện ổ khóa ở cột
HttpOnly và \code{document.cookie} trả về chuỗi rỗng.
\end{handson}

\section{Ngoài dự án}

Mẫu ở trên --- access token trong bộ nhớ, refresh token trong cookie
\code{HttpOnly}, refresh âm thầm khi tải --- là thứ SPA SDK của Auth0,
của Okta, và phần lớn ứng dụng web ngân hàng hiện thực, thường kèm xoay
vòng refresh và phát hiện tái sử dụng ở server. Phương án thay thế loại
trình duyệt khỏi việc xử lý token hoàn toàn là
\emph{backend-for-frontend}: một server nhỏ (hoặc chính gateway) giữ
token trong session phía server, và trình duyệt chỉ mang một cookie
session bình thường; Spring Cloud Gateway hỗ trợ điều này với filter
\code{TokenRelay} và đăng nhập OIDC. Người phỏng vấn hỏi ``bạn cất JWT
ở đâu trong trình duyệt?'' mong đợi ``không ở đâu mà script đọc được,
và refresh token không bao giờ trong URL'' theo sau là bảng đánh đổi ở
trên. ``Còn XSS thì sao?'' có một câu trả lời không hề về nơi cất: một
Content Security Policy ngăn script lạ chạy ngay từ đầu, thứ phòng thủ
mà mọi lựa chọn nơi cất vẫn phụ thuộc vào.

\section{Câu hỏi từ phỏng vấn thật}

\subsection*{Q: Bạn cất JWT ở đâu trong trình duyệt, và vì sao?}

Bắt đầu bằng tách hai token, vì chúng có thời gian sống và bán kính
thiệt hại khác nhau. Access token 15 phút có thể sống trong bộ nhớ
JavaScript: mất nó khi tải lại tốn một lần refresh âm thầm, và script
trộm được nó có mười lăm phút. Refresh token là thông tin xác thực
quan trọng --- bảy ngày ở đây --- và nó thuộc về một cookie
\code{HttpOnly}, \code{Secure}, \code{SameSite=Strict} giới hạn bằng
\code{Path} tới endpoint refresh, để không script nào đọc được và trình
duyệt chỉ gửi nó tới một URL. \code{localStorage}, thứ dự án dùng hôm
nay, bào chữa được cho phiên bản đầu và không bào chữa được khi bất kỳ
script bên thứ ba nào chạy trên trang: cách một \code{JSON.parse} là
một tuần truy cập. Câu kết senior là nơi cất là tuyến thứ hai; tuyến
đầu là Content Security Policy ngăn script bị chèn thực thi ngay từ
đầu.

\subsection*{Q: Mười tab nhận 401 cùng một khoảnh khắc. Bao nhiêu lời gọi refresh xảy ra?}

Với interceptor như đã viết, mười --- mỗi tab có bộ nhớ riêng và
request đang chờ riêng, và trong một tab mỗi truy vấn song song chạy
interceptor riêng rẽ, nên có thể hơn mười. Với xoay vòng refresh token
ở server đó không phải lãng phí, đó là sự cố: lần refresh đầu xoay
token, chín cái kia đưa ra token cũ và bị từ chối, và cơ chế phát hiện
tái sử dụng mà server tốt hiện thực thu hồi cả họ token. Trong một tab
cách sửa là promise dùng chung --- 401 đầu khởi động refresh, các 401
sau await nó, tất cả phát lại với kết quả. Xuyên tab là
\code{navigator.locks}: các tab lấy khóa lần lượt, và tab thứ hai thấy
cookie đã xoay và thành công với một lần refresh bình thường. Điều đáng
chủ động nói thêm là server nên khoan dung một cửa sổ ân hạn ngắn trên
token trước đó, vài giây, vì việc mạng đảo thứ tự nghĩa là token ``cũ''
có thể tới hợp lệ sau lần xoay.

\subsection*{Q: XSS so với CSRF --- mỗi lựa chọn nơi cất phơi bạn ra trước cái nào?}

Bearer token trong storage: miễn nhiễm CSRF, vì trình duyệt không bao
giờ tự đính header \code{Authorization}, và phơi nhiễm hoàn toàn trước
XSS, vì bất kỳ script nào trên origin đọc được storage. Cookie: miễn
nhiễm việc XSS đọc khi \code{HttpOnly}, và phơi nhiễm trước CSRF vì
trình duyệt đính cookie vào bất kỳ request nào tới host đó, kể cả từ
trang của kẻ tấn công --- trừ khi \code{SameSite=Strict} hay
\code{Lax} chặn nó, và một kiểm tra header tùy biến hay token hỗ trợ
cho client cũ. Vậy lựa chọn thật là giữa một bài toán có cách sửa cơ
học (CSRF) và một bài toán không có (XSS). Kỹ sư senior thêm rằng XSS
trên trang có cookie \code{HttpOnly} vẫn có thể \emph{dùng} phiên trong
lúc trang mở --- chỉ là không rút được thông tin xác thực ra --- và đó
là lý do access token ngắn và thu hồi phía server vẫn quan trọng.

\subsection*{Q: Người dùng đóng laptop tám ngày rồi quay lại. Chuyện gì xảy ra?}

Cả hai thông tin xác thực đều chết: access token từ mười lăm phút,
refresh token từ một ngày. Lời gọi API đầu nhận 401, interceptor gọi
refresh, auth-service thấy dòng hết hạn, xóa nó và trả 401, client
đăng xuất và chuyển hướng tới \code{/login}. Đó là kết quả đúng; câu
hỏi chất lượng là người dùng thấy gì trên đường đi. Mọi truy vấn đang
bay thất bại trước, và nếu query client thử lại request thất bại ba lần
với backoff --- mặc định của TanStack Query --- chuyển hướng đến sau vài
giây vòng xoay và các 401 lặp lại trong log. Cấu hình query client
không thử lại với 401, chạy một lần refresh khi ứng dụng mount trước
khi render route được bảo vệ, và cú quay lại sau tám ngày là một lần
bật sạch về đăng nhập với route dự định được giữ cho sau đó.

\subsection*{Q: Làm sao đăng xuất một người dùng ở mọi nơi?}

Phía refresh là một thao tác cơ sở dữ liệu: xóa mọi dòng trong
\code{refresh\_tokens} của người dùng đó, không chỉ dòng trình duyệt
hiện tại đưa ra. Mọi thiết bị khác thất bại ở lần refresh kế tiếp và
về trang đăng nhập trong vòng mười lăm phút. Phía access là bài toán
của Chương~28 --- token đã ký không thể thu hồi --- nên hãy nói thật về
cửa sổ và nêu hai cách đóng nó: một deny-list theo \code{jti} trong
Redis được gateway kiểm tra, hoặc access token ngắn hơn. Trong dự án
này thay đổi cụ thể là một phương thức repository,
\code{deleteAllByUser(user)}, gọi từ một endpoint
\code{/auth/logout-all} mới và từ hành động vô hiệu người dùng của
admin. Nhắc cái bẫy phỏng vấn: một ``đăng xuất'' phía client chỉ dọn
storage không thu hồi gì cả --- token vẫn hợp lệ ở bất kỳ ai đã sao
chép nó.

\subsection*{Q: Lỗi CORS xuất hiện ở production nhưng không bao giờ ở development. Bạn kiểm tra gì?}

Trước hết, request có cross-origin ở production mà không ở development
hay không --- trong dự án này proxy của Vite và Traefik đều giữ trang
và API trên một origin, nên lỗi CORS ở production nghĩa là việc triển
khai đã tách chúng: CDN cho bundle, host hay cổng khác cho gateway. Rồi
so header \code{Origin} trình duyệt gửi, từng byte kể cả scheme và
cổng, với danh sách cho phép của gateway; \code{https://app.example}
và \code{https://app.example:443} là hai chuỗi khác nhau với Spring.
Rồi quy tắc credentials: với cookie hay \code{withCredentials},
origin wildcard bị trình duyệt từ chối ngay cả khi server gửi nó, và
\code{Access-Control-Allow-Credentials: true} phải có mặt. Rồi
preflight: một \code{OPTIONS} rơi vào route bỏ tiền tố và chuyển tiếp
tới service sẽ trả 401 của service đó thay vì các header CORS, nên CORS
phải được trả lời tại gateway trước khi định tuyến --- điều
\code{CorsWebFilter} làm. Cuối cùng, nhìn log server và mong đợi không
thấy gì: request ở đó thành công; trình duyệt mới là bên từ chối, và đó
là hình dạng của mọi bài toán CORS.

\begin{quiz}
\begin{enumerate}
  \item Mở DevTools trên ứng dụng đang chạy và nêu tên khóa storage giữ
  cả hai token. Viết biểu thức JavaScript một dòng mà một script bị
  chèn sẽ dùng để đọc refresh token, và nói nó làm được gì với token
  đó trong bảy ngày tới.
  \item Interceptor gọi \code{axios} trần cho lần refresh thay vì
  \code{apiClient}. Giải thích vòng lặp sẽ xảy ra nếu không, từng bước,
  bắt đầu từ một 401 trên \code{/refresh}.
  \item Với xoay vòng refresh token bật ở server, lần theo bốn 401 song
  song qua interceptor như đã viết và cho biết ở bước nào người dùng bị
  đăng xuất. Rồi giải thích vì sao phiên bản single-flight tránh được.
  \item Thiết kế cookie đặt \code{Path=/api/auth/refresh}. Nêu hai đòn
  tấn công khác nhau mà thuộc tính này làm cùn, và giải thích vì sao
  path được viết theo góc nhìn của gateway chứ không phải auth-service.
  \item \code{CorsConfig} liệt kê hai origin và cho phép credentials.
  Trong mode triển khai nào trong hai mode của dự án nó được hỏi tới,
  và một thay đổi triển khai duy nhất nào sẽ khiến nó quan trọng?
  \item Người dùng báo lỗi CORS; log gateway cho thấy request trả 200.
  Giải thích vì sao cả hai quan sát đúng cùng lúc, và liệt kê ba giá
  trị header bạn sẽ so sánh trước tiên.
\end{enumerate}
\end{quiz}
